% GENERATED FILE. Edit canonical lessons, then run npm run generate:books.
% canonical-source-hash: fnv1a64:077b4d2124079962
% canonical-lessons: RU-C157-goluboy, RU-C157-sladkiy, RU-C157-kislyy, RU-C157-vkusnyy, RU-C157-udobnyy, RU-R157-first-pass-words-from-chapters-149-152, RU-R157-second-pass-words-from-chapters-153-157

\chapter{Light blue, Sweet, Sour, Tasty, Comfortable}
\label{ch:ru-light-blue-sweet-sour-tasty-comfortable}
\hlchaptermodality{157}

\begin{chapteropening}
\textbf{By the end of this chapter:} \emph{I can say light blue, sweet, sour, tasty and comfortable.}

Complete the last lesson of chapter 157: I can say light blue, sweet, sour, tasty and comfortable.
\end{chapteropening}

\section[\texorpdfstring{\ru{голубой} (golubóy)}{голубой (golubóy)}]{\ru{голубой} (golubóy) — light blue}

\label{lesson:RU-C157-goluboy}



\begin{quote}
Before the new one: say the Russian for dirty, then the Russian for empty.
\end{quote}

\subsection*{You'll want to know: \ru{голубой}}
\textbf{\ru{голубой}} (\emph{golubóy}) — \textquotedblleft{}light blue\textquotedblright{}.

\begin{grammarlens}[title={What you've built}]
One more word for this chapter.
\end{grammarlens}

\subsection*{Your turn}
\begin{itemize}
\raggedright
  \item \emph{Say it:} \emph{golubóy}
  \item \emph{Say it:} \emph{golubóy}, once more
  \item \emph{Say it:} say \emph{golubóy}
  \item \emph{Recall:} say the Russian for dirty, then the Russian for empty, then say \emph{golubóy} again
\end{itemize}

\begin{tcolorbox}[breakable,colback=teal!4,colframe=teal!35!black,title={Before you move on}]
What does \ru{голубой} mean? (\textquotedblleft{}Light blue\textquotedblright{}.) Say it once more.
\end{tcolorbox}
\section[\texorpdfstring{\ru{сладкий} (sládkiy)}{сладкий (sládkiy)}]{\ru{сладкий} (sládkiy) — sweet (taste)}

\label{lesson:RU-C157-sladkiy}



\begin{quote}
Before the new one: say the Russian for empty, then the Russian for light blue.
\end{quote}

\subsection*{You'll want to know: \ru{сладкий}}
\textbf{\ru{сладкий}} (\emph{sládkiy}) — \textquotedblleft{}sweet (taste)\textquotedblright{}.

\begin{grammarlens}[title={What you've built}]
One more word for this chapter.
\end{grammarlens}

\subsection*{Your turn}
\begin{itemize}
\raggedright
  \item \emph{Say it:} \emph{sládkiy}
  \item \emph{Say it:} \emph{sládkiy}, once more
  \item \emph{Say it:} say \emph{sládkiy}
  \item \emph{Recall:} say the Russian for empty, then the Russian for light blue, then say \emph{sládkiy} again
\end{itemize}

\begin{tcolorbox}[breakable,colback=teal!4,colframe=teal!35!black,title={Before you move on}]
What does \ru{сладкий} mean? (\textquotedblleft{}Sweet (taste)\textquotedblright{}.) Say it once more.
\end{tcolorbox}
\section[\texorpdfstring{\ru{кислый} (kíslyy)}{кислый (kíslyy)}]{\ru{кислый} (kíslyy) — sour}

\label{lesson:RU-C157-kislyy}



\begin{quote}
Before the new one: say the Russian for light blue, then the Russian for sweet.
\end{quote}

\subsection*{You'll want to know: \ru{кислый}}
\textbf{\ru{кислый}} (\emph{kíslyy}) — \textquotedblleft{}sour\textquotedblright{}.

\begin{grammarlens}[title={What you've built}]
One more word for this chapter.
\end{grammarlens}

\subsection*{Your turn}
\begin{itemize}
\raggedright
  \item \emph{Say it:} \emph{kíslyy}
  \item \emph{Say it:} \emph{kíslyy}, once more
  \item \emph{Say it:} say \emph{kíslyy}
  \item \emph{Recall:} say the Russian for light blue, then the Russian for sweet, then say \emph{kíslyy} again
\end{itemize}

\begin{tcolorbox}[breakable,colback=teal!4,colframe=teal!35!black,title={Before you move on}]
What does \ru{кислый} mean? (\textquotedblleft{}Sour\textquotedblright{}.) Say it once more.
\end{tcolorbox}
\section[\texorpdfstring{\ru{вкусный} (vkúsnyy)}{вкусный (vkúsnyy)}]{\ru{вкусный} (vkúsnyy) — tasty}

\label{lesson:RU-C157-vkusnyy}



\begin{quote}
Before the new one: say the Russian for sweet, then the Russian for sour.
\end{quote}

\subsection*{You'll want to know: \ru{вкусный}}
\textbf{\ru{вкусный}} (\emph{vkúsnyy}) — \textquotedblleft{}tasty\textquotedblright{}.

\begin{grammarlens}[title={What you've built}]
One more word for this chapter.
\end{grammarlens}

\subsection*{Your turn}
\begin{itemize}
\raggedright
  \item \emph{Say it:} \emph{vkúsnyy}
  \item \emph{Say it:} \emph{vkúsnyy}, once more
  \item \emph{Say it:} say \emph{vkúsnyy}
  \item \emph{Recall:} say the Russian for sweet, then the Russian for sour, then say \emph{vkúsnyy} again
\end{itemize}

\begin{tcolorbox}[breakable,colback=teal!4,colframe=teal!35!black,title={Before you move on}]
What does \ru{вкусный} mean? (\textquotedblleft{}Tasty\textquotedblright{}.) Say it once more.
\end{tcolorbox}
\section[\texorpdfstring{\ru{удобный} (udóbnyy)}{удобный (udóbnyy)}]{\ru{удобный} (udóbnyy) — comfortable, convenient}

\label{lesson:RU-C157-udobnyy}



\begin{quote}
Before the new one: say the Russian for sour, then the Russian for tasty.
\end{quote}

\subsection*{You'll want to know: \ru{удобный}}
\textbf{\ru{удобный}} (\emph{udóbnyy}) — \textquotedblleft{}comfortable, convenient\textquotedblright{}.

\begin{grammarlens}[title={What you've built}]
One more word for this chapter.
\end{grammarlens}

\subsection*{Your turn}
\begin{itemize}
\raggedright
  \item \emph{Say it:} \emph{udóbnyy}
  \item \emph{Say it:} \emph{udóbnyy}, once more
  \item \emph{Say it:} say \emph{udóbnyy}
  \item \emph{Recall:} say the Russian for sour, then the Russian for tasty, then say \emph{udóbnyy} again
\end{itemize}

\begin{tcolorbox}[breakable,colback=teal!4,colframe=teal!35!black,title={Before you move on}]
What does \ru{удобный} mean? (\textquotedblleft{}Comfortable, convenient\textquotedblright{}.) Say it once more.
\end{tcolorbox}
\section[(dialogue)]{Words from chapters 149-152 — a first pass}

\label{lesson:RU-R157-first-pass-words-from-chapters-149-152}



\begin{quote}
Say the words for tasty and comfortable.
\end{quote}

\subsection*{Your turn}
Say these aloud:

\begin{itemize}
\raggedright
  \item half an hour, a quarter, a day and a night, dawn, sunset
  \item rarely, every, at first, at last, still; more
  \item a place; a seat, the centre, the outskirts, a district, an avenue
  \item beside; roughly, around, inside, outside, behind
\end{itemize}

\begin{tcolorbox}[breakable,colback=teal!4,colframe=teal!35!black,title={Before you move on}]
Half an hour? (\textbf{\ru{полчаса}}.) A quarter? (\textbf{\ru{четверть}}.) A day and a night? (\textbf{\ru{сутки}}.) Dawn? (\textbf{\ru{рассвет}}.) Sunset? (\textbf{\ru{закат}}.) Rarely? (\textbf{\ru{редко}}.) Every? (\textbf{\ru{каждый}}.) At first? (\textbf{\ru{сначала}}.) At last? (\textbf{\ru{наконец}}.) Still; more? (\textbf{\ru{ещё}}.) A place; a seat? (\textbf{\ru{место}}.) The centre? (\textbf{\ru{центр}}.) The outskirts? (\textbf{\ru{окраина}}.) A district? (\textbf{\ru{район}}.) An avenue? (\textbf{\ru{проспект}}.) Beside; roughly? (\textbf{\ru{около}}.) Around? (\textbf{\ru{вокруг}}.) Inside? (\textbf{\ru{внутри}}.) Outside? (\textbf{\ru{снаружи}}.) Behind? (\textbf{\ru{сзади}}.)
\end{tcolorbox}
\section[(dialogue)]{Words from chapters 153-157 — a second pass}

\label{lesson:RU-R157-second-pass-words-from-chapters-153-157}



\begin{quote}
Say the words for a gym and a mosque.
\end{quote}

\subsection*{Your turn}
Say these aloud:

\begin{itemize}
\raggedright
  \item a gym, a mosque, a temple, a monument, a fountain
  \item so, since, if, although, otherwise
  \item nice, clever, silly, lazy, polite
  \item rich, poor, clean, dirty, empty
  \item light blue, sweet, sour, tasty, comfortable
\end{itemize}

\begin{tcolorbox}[breakable,colback=teal!4,colframe=teal!35!black,title={Before you move on}]
A gym? (\textbf{\ru{спортзал}}.) A mosque? (\textbf{\ru{мечеть}}.) A temple? (\textbf{\ru{храм}}.) A monument? (\textbf{\ru{памятник}}.) A fountain? (\textbf{\ru{фонтан}}.) So? (\textbf{\ru{поэтому}}.) Since? (\textbf{\ru{так} \ru{как}}.) If? (\textbf{\ru{если}}.) Although? (\textbf{\ru{хотя}}.) Otherwise? (\textbf{\ru{иначе}}.) Nice? (\textbf{\ru{милый}}.) Clever? (\textbf{\ru{умный}}.) Silly? (\textbf{\ru{глупый}}.) Lazy? (\textbf{\ru{ленивый}}.) Polite? (\textbf{\ru{вежливый}}.) Rich? (\textbf{\ru{богатый}}.) Poor? (\textbf{\ru{бедный}}.) Clean? (\textbf{\ru{чистый}}.) Dirty? (\textbf{\ru{грязный}}.) Empty? (\textbf{\ru{пустой}}.) Light blue? (\textbf{\ru{голубой}}.) Sweet? (\textbf{\ru{сладкий}}.) Sour? (\textbf{\ru{кислый}}.) Tasty? (\textbf{\ru{вкусный}}.) Comfortable? (\textbf{\ru{удобный}}.)
\end{tcolorbox}
